% GENERATED FILE. Edit canonical lessons, then run npm run generate:books.
% canonical-source-hash: fnv1a64:1e55d5909c356091
% canonical-lessons: GU-C410-bathrum, GU-C410-sandas, GU-C410-gijhar, GU-C410-ulecvu, GU-C410-nicovvu

\chapter{Bathrooms and Hot Water}
\label{ch:gu-bathrooms-and-hot-water}
\hlchaptermodality{410}

\begin{chapteropening}
\textbf{By the end of this chapter:} \emph{I can ask about the bathroom and the hot water.}

Complete the last lesson of chapter 410: I can ask about the bathroom and the hot water.
\end{chapteropening}

\section[\texorpdfstring{\gu{બાથરૂમ} (bāthrūm)}{બાથરૂમ (bāthrūm)}]{\gu{બાથરૂમ} (bāthrūm) — a bathroom}

\label{lesson:GU-C410-bathrum}



\begin{quote}
Before the new one: say the Gujarati for an inconvenience, then the Gujarati for to welcome.
\end{quote}

\subsection*{You'll want to know: \gu{બાથરૂમ}}
\textbf{\gu{બાથરૂમ}} — \emph{bāthrūm} — \textquotedblleft{}a bathroom\textquotedblright{}.

\begin{quote}
\gu{અમારા} \gu{ઓરડામાં} \gu{બાથરૂમ} \gu{નથી}.
\end{quote}

\begin{quote}
\emph{amārā orḍāmā̃ bāthrūm nathī.}
\end{quote}

\begin{quote}
There is no bathroom in our room.
\end{quote}

A neuter noun.

\begin{cousinweb}
\textbf{\gu{બાથરૂમ}} (\emph{bāthrūm}) is borrowed from English \emph{bathroom}.
\end{cousinweb}

\begin{grammarlens}[title={What you've built}]
One more word for getting around.
\end{grammarlens}

\subsection*{Your turn}
\begin{itemize}
\raggedright
  \item \emph{Say it:} \emph{bāthrūm}
  \item \emph{Say it:} \emph{bāthrūm}, once more
  \item \emph{Say it:} \emph{amārā orḍāmā̃ bāthrūm nathī.} — There is no bathroom in our room.
  \item \emph{Recall:} say the Gujarati for an inconvenience, then the Gujarati for to welcome, then say \emph{bāthrūm} again
\end{itemize}

\begin{tcolorbox}[breakable,colback=teal!4,colframe=teal!35!black,title={Before you move on}]
What does \gu{બાથરૂમ} mean? (\textquotedblleft{}A bathroom\textquotedblright{}.) Say it once more.
\end{tcolorbox}
\section[\texorpdfstring{\gu{સંડાસ} (sanḍās)}{સંડાસ (sanḍās)}]{\gu{સંડાસ} (sanḍās) — a toilet; a lavatory}

\label{lesson:GU-C410-sandas}



\begin{quote}
Before the new one: say the Gujarati for to welcome, then the Gujarati for a bathroom.
\end{quote}

\subsection*{You'll want to know: \gu{સંડાસ}}
\textbf{\gu{સંડાસ}} — \emph{sanḍās} — \textquotedblleft{}a toilet; a lavatory\textquotedblright{}.

\begin{quote}
\gu{સંડાસ} \gu{ક્યાં} \gu{છે}?
\end{quote}

\begin{quote}
\emph{sanḍās kyā̃ che?}
\end{quote}

\begin{quote}
Where is the toilet?
\end{quote}

A neuter noun.

\begin{cousinweb}
\textbf{\gu{સંડાસ}} (\emph{sanḍās}) No dictionary we use records an older form, so its history is not settled.
\end{cousinweb}

\begin{grammarlens}[title={What you've built}]
One more word for getting around.
\end{grammarlens}

\subsection*{Your turn}
\begin{itemize}
\raggedright
  \item \emph{Say it:} \emph{sanḍās}
  \item \emph{Say it:} \emph{sanḍās}, once more
  \item \emph{Say it:} \emph{sanḍās kyā̃ che?} — Where is the toilet?
  \item \emph{Recall:} say the Gujarati for to welcome, then the Gujarati for a bathroom, then say \emph{sanḍās} again
\end{itemize}

\begin{tcolorbox}[breakable,colback=teal!4,colframe=teal!35!black,title={Before you move on}]
What does \gu{સંડાસ} mean? (\textquotedblleft{}A toilet; a lavatory\textquotedblright{}.) Say it once more.
\end{tcolorbox}
\section[\texorpdfstring{\gu{ગીઝર} (gījhar)}{ગીઝર (gījhar)}]{\gu{ગીઝર} (gījhar) — a water heater; a geyser}

\label{lesson:GU-C410-gijhar}



\begin{quote}
Before the new one: say the Gujarati for a bathroom, then the Gujarati for a toilet.
\end{quote}

\subsection*{You'll want to know: \gu{ગીઝર}}
\textbf{\gu{ગીઝર}} — \emph{gījhar} — \textquotedblleft{}a water heater; a geyser\textquotedblright{}.

\begin{quote}
\gu{હોટેલના} \gu{ગીઝરમાં} \gu{ગરમ} \gu{પાણી} \gu{નથી}.
\end{quote}

\begin{quote}
\emph{hoṭelnā gījharmā̃ garam pāṇī nathī.}
\end{quote}

\begin{quote}
There is no hot water in the hotel's water heater.
\end{quote}

A neuter noun.

\begin{cousinweb}
\textbf{\gu{ગીઝર}} (\emph{gījhar}) is borrowed from English \emph{geyser}, which in India means a water heater in the bathroom.
\end{cousinweb}

\begin{grammarlens}[title={What you've built}]
One more word for getting around.
\end{grammarlens}

\subsection*{Your turn}
\begin{itemize}
\raggedright
  \item \emph{Say it:} \emph{gījhar}
  \item \emph{Say it:} \emph{gījhar}, once more
  \item \emph{Say it:} \emph{hoṭelnā gījharmā̃ garam pāṇī nathī.} — There is no hot water in the hotel's water heater.
  \item \emph{Recall:} say the Gujarati for a bathroom, then the Gujarati for a toilet, then say \emph{gījhar} again
\end{itemize}

\begin{tcolorbox}[breakable,colback=teal!4,colframe=teal!35!black,title={Before you move on}]
What does \gu{ગીઝર} mean? (\textquotedblleft{}A water heater; a geyser\textquotedblright{}.) Say it once more.
\end{tcolorbox}
\section[\texorpdfstring{\gu{ઉલેચવું} (ulecvũ)}{ઉલેચવું (ulecvũ)}]{\gu{ઉલેચવું} (ulecvũ) — to bail out (water)}

\label{lesson:GU-C410-ulecvu}



\begin{quote}
Before the new one: say the Gujarati for a toilet, then the Gujarati for a water heater.
\end{quote}

\subsection*{You'll want to know: \gu{ઉલેચવું}}
\textbf{\gu{ઉલેચવું}} — \emph{ulecvũ} — \textquotedblleft{}to bail out (water)\textquotedblright{}.

\begin{quote}
\gu{વરસાદ} \gu{પછી} \gu{અમે} \gu{ઓરડામાંથી} \gu{પાણી} \gu{ઉલેચ્યું}.
\end{quote}

\begin{quote}
\emph{varsād pachī ame orḍāmā̃thī pāṇī ulecyũ.}
\end{quote}

\begin{quote}
After the rain we bailed the water out of the room.
\end{quote}

Past, with \gu{તેણે} (he or she …): \gu{તેણે} \gu{પાણી} \gu{ઉલેચ્યું}, \emph{teṇe pāṇī ulecyũ}.

\begin{cousinweb}
\textbf{\gu{ઉલેચવું}} (\emph{ulecvũ}) goes back to \emph{*udriñcati}, \textquotedblleft{}empties\textquotedblright{}, a form CDIAL lists beside Sanskrit \emph{údricyatē}; CDIAL lists the Gujarati verb, \textquotedblleft{}to draw out water\textquotedblright{} (CDIAL 2061).
\end{cousinweb}

\begin{grammarlens}[title={What you've built}]
One more word for getting around.
\end{grammarlens}

\subsection*{Your turn}
\begin{itemize}
\raggedright
  \item \emph{Say it:} \emph{ulecvũ}
  \item \emph{Say it:} \emph{ulecvũ}, once more
  \item \emph{Say it:} \emph{varsād pachī ame orḍāmā̃thī pāṇī ulecyũ.} — After the rain we bailed the water out of the room.
  \item \emph{Recall:} say the Gujarati for a toilet, then the Gujarati for a water heater, then say \emph{ulecvũ} again
\end{itemize}

\begin{tcolorbox}[breakable,colback=teal!4,colframe=teal!35!black,title={Before you move on}]
What does \gu{ઉલેચવું} mean? (\textquotedblleft{}To bail out (water)\textquotedblright{}.) Say it once more.
\end{tcolorbox}
\section[\texorpdfstring{\gu{નિચોવવું} (nicovvũ)}{નિચોવવું (nicovvũ)}]{\gu{નિચોવવું} (nicovvũ) — to wring out; to squeeze out (water)}

\label{lesson:GU-C410-nicovvu}



\begin{quote}
Before the new one: say the Gujarati for a water heater, then the Gujarati for to bail out.
\end{quote}

\subsection*{You'll want to know: \gu{નિચોવવું}}
\textbf{\gu{નિચોવવું}} — \emph{nicovvũ} — \textquotedblleft{}to wring out; to squeeze out (water)\textquotedblright{}.

\begin{quote}
\gu{મીરાએ} \gu{નાહીને} \gu{ભીનો} \gu{ટુવાલ} \gu{નિચોવ્યો}.
\end{quote}

\begin{quote}
\emph{Mīrāe nāhīne bhīno ṭuvāl nicovyo.}
\end{quote}

\begin{quote}
After her bath Mira wrung out the wet towel.
\end{quote}

Past, with \gu{તેણે} (he or she …): \gu{તેણે} \gu{ટુવાલ} \gu{નિચોવ્યો}, \emph{teṇe ṭuvāl nicovyo}.

\begin{cousinweb}
\textbf{\gu{નિચોવવું}} (\emph{nicovvũ}) goes back to \emph{*niścyōtayati}, \textquotedblleft{}causes to trickle out\textquotedblright{}, a form CDIAL reconstructs; CDIAL lists the Gujarati verb, \textquotedblleft{}to squeeze out\textquotedblright{} (CDIAL 7451): you twist a cloth until the water runs out.
\end{cousinweb}

\begin{grammarlens}[title={What you've built}]
One more word for getting around.
\end{grammarlens}

\subsection*{Your turn}
\begin{itemize}
\raggedright
  \item \emph{Say it:} \emph{nicovvũ}
  \item \emph{Say it:} \emph{nicovvũ}, once more
  \item \emph{Say it:} \emph{Mīrāe nāhīne bhīno ṭuvāl nicovyo.} — After her bath Mira wrung out the wet towel.
  \item \emph{Recall:} say the Gujarati for a water heater, then the Gujarati for to bail out, then say \emph{nicovvũ} again
\end{itemize}

\begin{tcolorbox}[breakable,colback=teal!4,colframe=teal!35!black,title={Before you move on}]
What does \gu{નિચોવવું} mean? (\textquotedblleft{}To wring out; to squeeze out (water)\textquotedblright{}.) Say it once more.
\end{tcolorbox}
